%%%%%%%%%%%%%%%%%%%%%%%%%%%%%%%%%%%%%%%%%%%%%%%%%%%%%%%%%%%%%%%%%%%%%%%%%%%%%%%%%%%%%%%%%%
% Inferences
%%%%%%%%%%%%%%%%%%%%%%%%%%%%%%%%%%%%%%%%%%%%%%%%%%%%%%%%%%%%%%%%%%%%%%%%%%%%%%%%%%%%%%%%%%

{\it In this section, we discuss an approach to fitting the data when there is a reasonable
number of events that survive. This is an approach that we could use as a cross-check
with a relaxed selection criterial on the classifier output, or if we want to attempt to 
measure the baryon-number conserving decay $B \rightarrow \bar{p} \Lambda^0$. An approach 
to extracting an upper limit or the significance on an observed signal will be discussed
in a later section.}\\

